\subsubsection{2.1 研究意义}

当前，人工智能、数据要素、边缘计算与电磁频谱感知技术加速融合，电磁频谱数据智能处理正在成为复杂环境目标感知侦测的重要支撑。十五五时期，国家围绕数字中国建设、人工智能+行动和海洋强国建设作出系统部署。《中华人民共和国国民经济和社会发展第十五个五年规划纲要》提出推进数字中国建设、促进人工智能创新应用、拓展海洋经济发展空间。《中共中央关于制定国民经济和社会发展第十五个五年规划的建议》强调全面实施人工智能+行动，以人工智能引领科研范式变革，并提出坚持陆海统筹，提高经略海洋能力，推动海洋经济高质量发展，加快建设海洋强国。《国务院关于深入实施人工智能+行动的意见》提出推动人工智能与经济社会各行业各领域广泛深度融合。《数字中国建设整体布局规划》提出夯实数字基础设施和数据资源体系。在上述政策牵引下，复杂海洋环境下电磁频谱数据的智能获取、有效存储和即时处理能力，正成为支撑海洋强国建设、数字中国建设和人工智能+行动落地的重要基础。

随着海上无人目标、无AIS目标、异常无人艇、高速小艇和非法作业船艇等活动场景日益复杂，复杂海况、雨雾夜间、海面起伏、多径传播和海杂波干扰对电磁感知侦测能力提出了更高要求。此类目标携带的通信、导航、遥控等电磁辐射源信号往往能量弱、持续时间短、时频特征不稳定，易被海杂波背景掩盖；同时，海杂波场景电磁频谱数据连续产生、体量大、维度高，有限算力端侧难以对全量数据进行长期存储和即时处理，导致海杂波影响下无人目标电磁感知侦测性能受到严重限制。当前技术路径面临的突出矛盾在于，复杂海洋环境中微弱电磁频谱信号难以稳定提取，而端侧算力、存储和链路资源又难以支撑全量频谱数据的即时智能处理。在此背景下，发展面向海杂波场景的电磁频谱数据智能处理机理与方法，成为提升无人目标电磁感知侦测能力、支撑复杂海洋环境下频谱数据有效利用的重要方向。

当前，海杂波场景电磁频谱数据智能处理面临三方面关键挑战：

\textbf{频谱提取难度高。}在复杂海洋电磁环境中，海杂波呈现非高斯、非线性以及时空多径耦合特征，海上目标产生的微弱辐射信号信噪比极低，且多表现为短时猝发形态，极易被强杂波完全遮蔽，实现目标频谱有效辨识成为首要技术壁垒。传统恒虚警检测、空时自适应处理算法均建立在杂波统计平稳的假设之上，面对动态变化的海杂波会出现检测门限偏移、杂波对消残留过多等问题，难以高保真捕获瞬态微弱信号。一方面，动态海况下海杂波与目标信号的耦合传播规律难以精准刻画，缺乏完善的时空耦合机理表征模型；另一方面，低信噪比短时信号难以获取足量标注样本，常规信号分离算法容易丢失目标关键特征。仅依靠传统滤波手段无法兼顾检出率与虚警率，融合物理传播规律与小样本学习的一体化提取框架仍未成熟。因此，构建适配非平稳海杂波的信号建模、干扰对消与特征增强体系，是实现微弱频谱精准提取必须突破的技术瓶颈。

\textbf{海量数据载荷大。}海域广域自主侦测过程中会持续生成海量高维频谱数据，给端侧设备存储、检索与运算带来巨大载荷，同时设备算力与硬件资源存在天然约束，如何实现数据高效管控成为核心难题。现有端侧存储方案多采用固定阈值过滤、均匀采样模式，无法对频谱数据开展细粒度价值评估，易丢失含目标的关键信息，粗放式压缩还会造成频谱语义失真。此外，海上侦测环境对抗性强，传统中心化存储模式存在数据被恶意篡改的风险，难以满足执法取证对数据溯源、完整性的硬性要求。端侧硬件无法运行大型评估模型与高开销共识加密算法，面向频谱场景的专用压缩算法、多级索引优化方案尚不健全，链上链下协同的轻量可信存储架构也未形成完整体系。在有限硬件条件下统筹数据筛选、压缩、索引与可信存证，目前尚无成熟解决方案，是频谱数据高效管理亟待攻克的技术难题。

\textbf{智能决策效率低。}海域自主侦测任务面临海杂波动态演变、辐射源类型繁杂、通信链路间歇中断等问题，单节点探测存在观测盲区，多节点协同又受带宽、节点信任问题制约，边缘端智能决策的时效性与可靠性难以保障。大参数量深度学习模型因算力、功耗限制无法部署于边缘设备，难以同时满足微弱信号识别高精度与响应高实时性的双重要求。现有多源数据融合方案大多依赖全量数据回传与云端集中决策，不仅无法适配海上弱连接通信环境，还易遭受被俘节点伪造数据的干扰。当前适配多海况、多辐射源的混合专家轻量化推理架构、知识驱动在线学习机制仍不完善，弱链路下自适应路由、稀疏传输策略以及分布式节点可信审计、抗抵赖流程也存在明显短板。在算力、时延、带宽多重约束下，搭建高精度、快响应、高可信的边缘智能协同处理框架，是保障广域海域自主侦测稳定运行必须跨越的技术鸿沟。

为了应对以上三大挑战，本项目拟围绕海杂波场景下电磁频谱数据智能处理机理与方法，系统构建从基础机理建模、算法方法设计到体系架构搭建的全链条研究路径。首先，以 “海杂波时空耦合传播解析” 为指导思想，开展微弱电磁频谱信号传播规律分析与精准提取研究：通过复杂海洋环境下频谱时空耦合传播机理建模、空域协同干扰对消实现目标频谱特征增强，融合时域、频域、空域多维联合表征完成强干扰下微弱信号高精度分离提取，为复杂海况环境下弱小目标频谱 “看得清” 奠定底层理论与算法基础；然后，以 “端侧轻量可信数据管控” 为指导思想，构建有限算力条件下频谱数据存储与可信管理体系：依托事件窗口搭建自监督频谱数据价值评估模型实现无效杂波前置过滤，设计频谱语义表征压缩与多级索引优化向量数据库存储载荷，采用链上链下协同架构实现频谱采集、存储、调用全流程可追溯防篡改管控，在资源约束下达成海量频谱数据 “存得、证得真”；接下来，以 “知识驱动边缘智能协同侦测” 为指导思想，搭建面向海域自主侦测的边缘频谱智能处理框架：构建海杂波 - 目标态势一体化频谱知识图谱并支持动态增量演化，结合知识引导轻量化模型完成微弱频谱识别与目标轨迹解算，设计多节点频谱认知协同机制与分布式可信决策流程，在弱带宽、低算力约束下实现侦测态势 “判得快、判得稳”；最后，集成本项目提出的传播模型、信号提取算法、轻量存储架构与边缘智能协同体系，搭建海杂波频谱智能处理一体化原型系统，开展多梯度海况场景仿真与海上外场实测闭环验证，打通频谱采集流、数据存储流、智能决策流，形成可落地、可迁移的海域侦测技术资产，为海杂波场景电磁频谱智能侦测体系提供完整理论支撑与工程示范。

本项目的成功实施有助于：(1) \textbf{突破海洋电磁频谱处理理论瓶颈，构建海杂波场景微弱信号与端边协同智能处理原创理论体系。}现有研究缺乏海杂波时空耦合传播、微弱频谱提取、端侧数据管控与边缘侦测协同的成套理论框架。本项目建立完整的机理建模、信号增强、数据管理与认知决策理论体系，填补海洋电磁智能侦测交叉学科空白，助力我国海洋感知基础研究跻身国际前列。
(2) \textbf{破解低算力端侧频谱存储与轻量化模型部署工程难题，支撑海域侦测装备规模化应用。}传统大模型、重型数据库算力开销高，无法适配海上小型侦测终端。本项目依托频谱压缩、自监督筛杂、轻量知识推理、链上链下可信存证技术搭建原型系统，为海上执法、航道监控、近海安防提供可复制技术方案，赋能海洋防控新质生产力发展。
(3) \textbf{提升海洋电磁装备自主可控水平，服务海洋强国与平安海岸建设。}面对近海安防压力与复杂海况自主作业刚需，国产化频谱侦测技术是海防产业链安全关键根基。项目成果可强化微弱目标探测、可信数据存储、分布式协同研判国产化能力，提升海域安防实力，为数字海洋与海洋领域高质量发展提供科技支撑。

综上所述，海杂波场景电磁频谱数据智能处理机理与方法研究，既是海洋电磁、信号处理、边缘 AI 多学科交叉演进的必然方向，也是化解海上安防、海域侦测等关键科技短板的现实需求，更是争夺海洋智能感知技术主动权、保障海洋科技安全的重要基石。开展本项目研究具备显著的必要性与研究价值。
